\appendix
\section{Comparison on noncompact domains and Fourier identification}
The main particle theorem is already complete. We record two analytic
details separately: a comparison theorem with precise noncompact
pressure assumptions, useful for other prescribed flows, and the
elementary Fourier uniqueness argument used in the periodic projection.

\subsection{The precise comparison class}\label{app:comparison}
Let $\Omega$ be $\T^3$, $\R^3$, or
$\R^2\times\T$. A classical solution on $[0,T]$ has continuous
$u$, $u_t$, all spatial derivatives of $u$ through order two, $p$, and
$\nabla p$ on the closed time cylinder, and satisfies the equation
pointwise. In the noncompact cases we additionally require
\begin{equation}\label{eq:energy-class}
 \begin{gathered}
 u\in C([0,T];H^2(\Omega))\cap C^1([0,T];L^2(\Omega)),
 \qquad p\in C([0,T];H^1(\Omega)),\\
 \sup_{[0,T]\times\Omega}(|u|+|\nabla u|)<\infty.
 \end{gathered}
\end{equation}
The Sobolev spaces are periodic in the torus coordinates. On $\T^3$
the pressure has mean zero; on the noncompact domains its $H^1$
membership fixes the representative, since a nonzero spatial constant
is not in $L^2$.

\begin{lemma}[Prescribed motion and uniqueness]\label{lem:realization}
Fix $\nu>0$ and $\Omega\in\{\T^3,\R^3,\R^2\times\T\}$.
Let $U$ be smooth and divergence free on $[0,\infty)\times\Omega$,
with $U(0)=0$, and suppose $U$ and $\nabla U$ are bounded on every
finite closed time cylinder. In the noncompact cases assume the
velocity conditions in \eqref{eq:energy-class}. Then the force
\[
 \mathcal F_\nu[U]=U_t+(U\cdot\nabla)U-\nu\Delta U
\]
has the solution $(u,p)=(U,0)$, unique in the classical class just
defined on each finite time interval. Its material trajectories exist
uniquely for every finite time. On $\T^3$, mean-zero $U$ gives
mean-zero $\mathcal F_\nu[U]$.

More generally, if an independently constructed pair $(U,P)$ is a
classical solution in the same class for a given force and initial
velocity, it is the unique solution for those data in that class.
In either noncompact uniqueness statement, the competitor's gradient
bound can be omitted while keeping its bounded velocity and all other
conditions. The reference velocity $U$ still has bounded gradient.
For smooth $U$, each finite-time material map $\Phi_t$ is a smooth
diffeomorphism and satisfies
\[
 U(t,\cdot)=(\partial_t\Phi_t)\circ\Phi_t^{-1},\qquad
 \partial_{tt}\Phi_t(a)
   =(f-\nabla P+\nu\Delta U)(t,\Phi_t(a)).
\]
\end{lemma}
\begin{proof}
Substitution proves the residual assertion. For either uniqueness
statement, let $(v,p)$ be another solution for the same data, and set
$w=v-U$ and $q=p-P$. Subtracting the equations gives
\[
 w_t+(v\cdot\nabla)w+(w\cdot\nabla)U
   =-\nabla q+\nu\Delta w,\qquad \diver w=0.
\]
On the torus, multiplication by $w$ and periodic integration give
\begin{equation}\label{eq:energy-comparison}
 \frac12\frac{d}{dt}\|w\|_2^2+\nu\|\nabla w\|_2^2
 \leq\|\nabla U(t)\|_{\infty,\mathrm{op}}\|w\|_2^2.
\end{equation}
The transport and pressure integrals vanish by incompressibility.

Here are the cutoff details on the other two domains. Let $x'$ denote
the noncompact coordinates, choose $0\leq\zeta\leq1$ smooth with
compact support, equal to one on $|x'|\leq1$, and set
$\zeta_R(x')=\zeta(x'/R)$. Multiply the difference equation by
$\zeta_Rw$ and integrate in all coordinates. Integration by parts
leaves, besides the term $-\int\zeta_Rw\cdot(w\cdot\nabla)U$,
three boundary errors. Their absolute values are at most
\[
 \frac C R\|v\|_\infty\|w\|_2^2,\qquad
 \frac C R\|q\|_2\|w\|_2,\qquad
 \frac {C\nu}R\|\nabla w\|_2\|w\|_2,
\]
respectively for transport, pressure, and viscosity. For example,
the pressure term is $\int q\nabla\zeta_R\cdot w$ because
$\diver w=0$. Each displayed norm is uniformly bounded on $[0,T]$
by \eqref{eq:energy-class}. Also $w\in C^1([0,T];L^2)$
justifies differentiating $\int\zeta_R|w|^2$. Integrate the identity
between any two times and then let $R\to\infty$. Dominated
convergence applies to the quadratic terms and the error integrals
tend to zero. This gives \eqref{eq:energy-comparison} in integrated
form, equivalently almost everywhere in time. Thus no decay of the
pressure beyond the stated $C_tH^1$ condition has been assumed.

On every finite interval the coefficient in
\eqref{eq:energy-comparison} is bounded. Its integrating factor and
$w(0)=0$ imply $w=0$. The equation then gives $\nabla q=0$;
the mean-zero or $H^1$ pressure convention gives $q=0$.
The estimates use boundedness of $v$ but never a bound on $\nabla v$,
which proves the stated competitor extension.

Finally, bounded spatial gradient gives a uniformly Lipschitz vector
field on each finite time interval, so the material ODE has unique
local solutions. Bounded velocity prevents escape to infinity in
finite time. Successive local solutions therefore cover each finite
interval, and periodic coordinates descend to their quotients.
Solving the same ODE backwards on $[0,t]$ gives the inverse map.
Smooth dependence on the initial point proves the diffeomorphism
assertion; differentiating the trajectory equation proves the two
displayed identities.
For the mean assertion, integrate the residual: the integrals of
$\Delta U$ and $(U\cdot\nabla)U=\diver(U\otimes U)$ vanish on
the torus.
\end{proof}


\subsection{Identifying continuous periodic functions}\label{app:fourier}
A continuous function on $\T^3$ is determined by its Fourier
coefficients. Here is the approximate-identity argument used in
Lemma~\ref{lem:projection}.

In one variable,
\[
 F_N(s)=\frac1N\left|\sum_{j=0}^{N-1}e^{2\pi i js}\right|^2
\]
is nonnegative, has integral one, and has Fourier coefficients
$(1-|k|/N)_+$. Outside any fixed neighborhood of the integers,
the geometric-sum formula bounds $F_N$ by $C/N$. Hence the product
of three such kernels has mass tending to one in every neighborhood
of zero. Uniform continuity shows that convolution with this product
converges uniformly to any continuous function on $\T^3$.
Vanishing Fourier coefficients therefore imply a vanishing function.
